\phantomsection
\section*{Введение}
\addcontentsline{toc}{chapter}{Введение}

Модели общего динамического равновесия занимают важное место в современной макроэкономической теории \cite{heermaussner2009dynamic}. Они позволяют исследовать совместную динамику выпуска, потребления, капитала и труда, а также анализировать механизм межвременного выбора экономических агентов. Одной из наиболее известных конструкций такого типа является модель реального делового цикла (Real Business Cycle, RBC), в которой колебания основных макроэкономических переменных объясняются реакцией экономики на реальные шоки при рациональном поведении домохозяйств и фирм \cite{kydland1982time}. Дальнейшим развитием этого подхода стали неокейнсианские DSGE-модели (Dynamic Stochastic General Equilibrium), которые дополнительно учитывают номинальные жёсткости (цен и заработных плат), несовершенную конкуренцию и денежно-кредитную политику \cite{smets2007shocks, christiano2005nominal}.

Модели реального делового цикла (RBC) и их неокейнсианские расширения (DSGE) являются стандартным инструментом макроэкономического анализа, используемым центральными банками, инвестиционными фондами и министерствами экономики для прогнозирования, стресс-тестирования и оценки последствий политических решений \cite{motto2025chapter, ristiniemi2025chapter}. Согласно опросу 22 центральных банков, DSGE-модели служат основной или вспомогательной моделью для большинства из них \cite{zamrazilova2025role}. Однако даже детерминированная RBC-модель приводит к нелинейной системе разностных уравнений, аналитическое решение которой невозможно. Традиционные локальные методы (линеаризация) дают приемлемую точность лишь в малой окрестности стационарного состояния \cite{judd1998numerical}, тогда как реальные шоки могут быть значительными. Поэтому актуальной задачей является реализация и сопоставление глобальных численных методов (например, метода стрельбы) с локальными подходами для построения переходных траекторий при больших отклонениях от равновесия \cite{lipton1983multiple}.

Целью выпускной квалификационной работы является сопоставление глобального (метод стрельбы) и локального (линеаризация) подходов к решению детерминированной модели реального делового цикла для закрытой экономики \cite{judd1998numerical}, а также анализ её стационарного состояния и переходной динамики \cite{heermaussner2009dynamic}. Дополнительно модель расширяется на стохастический случай для анализа импульсных откликов на шоки совокупной факторной производительности и денежно-кредитной политики \cite{sims2002solving, blanchardkahn1980}.

Объектом исследования является детерминированная и стохастическая макроэкономическая модель закрытой экономики типа RBC. Предметом исследования выступают стационарное состояние модели, переходная динамика капитала, потребления, выпуска и труда, а также численные методы построения устойчивой траектории \cite{judd1998numerical, fernandez2015solution}.

Работа состоит из введения, трёх глав, заключения и списка литературы. В первой главе формально ставятся задачи домохозяйства и фирмы, выводится необходимое условие экстремума, исследуется стационарное состояние и формируется итоговая система уравнений детерминированной модели в динамической форме. Во второй главе рассматриваются существующие численные подходы к решению RBC-моделей: даётся классификация на локальные и глобальные методы, обосновывается выбор методов для данной работы, излагается алгоритм построения траектории при фиксированном начальном уровне потребления, метод Ньютона для нахождения управляющей переменной и приводится псевдокод метода стрельбы. В третьей главе представлены результаты численных экспериментов: фазовые траектории в плоскостях $(K, C)$ и $(K, L)$, трёхмерная траектория, а также сравнение глобального и локального решений. В четвёртой главе формулируются основные выводы и обсуждаются перспективы дальнейших исследований.